\subsubsection{LCS / Edit distance}
\begin{lstlisting}[style=mystyle, language=C++]
// TC: O(N * M)
// usa: longest common subsequence y distancia de edicion (Levenshtein)
#include <bits/stdc++.h>
using namespace std;

int lcs(const string& a, const string& b){
	int n = a.size(), m = b.size();
	vector<vector<int>> dp(n + 1, vector<int>(m + 1, 0));
	for(int i=1;i<=n;i++){
		for(int j=1;j<=m;j++){
			if(a[i-1] == b[j-1]) dp[i][j] = dp[i-1][j-1] + 1;
			else dp[i][j] = max(dp[i-1][j], dp[i][j-1]);
		}
	}
	return dp[n][m];
}

int editDistance(const string& a, const string& b){
	int n = a.size(), m = b.size();
	vector<vector<int>> dp(n + 1, vector<int>(m + 1, 0));
	for(int i=0;i<=n;i++) dp[i][0] = i;
	for(int j=0;j<=m;j++) dp[0][j] = j;
	for(int i=1;i<=n;i++){
		for(int j=1;j<=m;j++){
			if(a[i-1] == b[j-1]) dp[i][j] = dp[i-1][j-1];
			else dp[i][j] = 1 + min({dp[i-1][j], dp[i][j-1], dp[i-1][j-1]});
		}
	}
	return dp[n][m];
}

int main(){
	cout << lcs("abcde", "ace") << "\n";       // 3
	cout << editDistance("kitten", "sitting") << "\n";  // 3
	return 0;
}
\end{lstlisting}
